\providecommand{\Res}[2]{\qty[round-mode=figures,round-precision=3]{\fpeval{#1}}{#2}}
\providecommand{\Val}[1]{\num[round-mode=figures,round-precision=4]{\fpeval{#1}}}
\providecommand{\ResSci}[2]{\qty[round-mode=figures,round-precision=3,exponent-mode=scientific]{\fpeval{#1}}{#2}}
\subsubsection*{Question a) Longueur d'onde}
\[ \lambda = \frac{c}{f} = \frac{\num{3e8}}{\num{\v{f}e6}} = \Res{3e8/\vsi{f}}{m} \]
La longueur d'onde est de \textbf{\Res{3e8/\vsi{f}}{m}}.

\subsubsection*{Question b) Premier maximum}
\begin{enumerate}[label=\arabic*)]
    \item \textbf{Variables connues et recherchées :}
    \begin{itemize}[label= -]
        \item $D = \qty{\v{D}}{m}$, $\lambda = \Res{3e8/\vsi{f}}{m}$, $m = 1$
        \item Variable recherchée : $\theta_1$
    \end{itemize}
    \item \textbf{Équations pertinentes :} sources en phase, maximum quand
    \[ D\sin\theta = m\lambda \]
    \item \textbf{Isolation de la variable recherchée :}
    \[ \theta_1 = \arcsin\left(\frac{\lambda}{D}\right) \]
    \item \textbf{Résolution numérique :}
    \[ \theta_1 = \arcsin\left(\frac{\Val{3e8/\vsi{f}}}{\num{\v{D}}}\right) = \Res{asind(3e8/(\vsi{f}*\v{D}))}{\degree} \]
    \item \textbf{Réponse :} le premier maximum est à \textbf{\Res{asind(3e8/(\vsi{f}*\v{D}))}{\degree}}.
\end{enumerate}

\subsubsection*{Question c) Nombre de directions de maximum}
\begin{enumerate}[label=\arabic*)]
    \item \textbf{Variables connues et recherchées :} $D/\lambda = \Res{\v{D}*\vsi{f}/3e8}{}$ ; variable recherchée : le nombre de valeurs entières de $m$
    \item \textbf{Équations pertinentes :} $\sin\theta = m\lambda/D$ doit respecter $|\sin\theta| \le 1$
    \item \textbf{Isolation de la variable recherchée :} $|m| \le D/\lambda$, donc $m_{\max} = \num{\fpeval{floor(\v{D}*\vsi{f}/3e8)}}$
    \item \textbf{Résolution numérique :} $m = -\num{\fpeval{floor(\v{D}*\vsi{f}/3e8)}}, \dots, 0, \dots, \num{\fpeval{floor(\v{D}*\vsi{f}/3e8)}}$, soit $2 \cdot \num{\fpeval{floor(\v{D}*\vsi{f}/3e8)}} + 1$ valeurs
    \item \textbf{Réponse :} il y a \textbf{\num{\fpeval{2*floor(\v{D}*\vsi{f}/3e8) + 1}}} directions de maximum.
\end{enumerate}

\subsubsection*{Question d) Premier minimum}
Minimum quand $D\sin\theta = (m + \tfrac{1}{2})\lambda$ ; le premier ($m = 0$) est à
\[ \theta = \arcsin\left(\frac{\lambda}{2D}\right) = \arcsin\left(\frac{\Val{3e8/\vsi{f}}}{2 \cdot \num{\v{D}}}\right) = \Res{asind(3e8/(2*\vsi{f}*\v{D}))}{\degree} \]
L'auditeur ne capte presque rien à \textbf{\Res{asind(3e8/(2*\vsi{f}*\v{D}))}{\degree}}.
